\begin{pastexam}[label={px:nsysu-115-dmalg-07}]{中山資工所資安組 115 學年度 第 8 題}
（設計、分析）

(Algorithm points) {[}10\%{]} Please describe the algorithm of Heap Sort and analyze its time complexity in detail with mathematical derivation.
\end{pastexam}

\begin{pastsoln}
\textbf{想法}。 最大堆積是一棵完全二元樹，存在陣列 $A[1..n]$：$A[i]$ 的子節點是 $A[2i]$、$A[2i + 1]$，每個節點都 $\ge$ 它的
子節點，所以最大值在 $A[1]$（資結書第~8.1~節）。 堆積排序（資結書第~14.6~節）分兩階段：先把整個陣列建成最大堆積；再反覆
把堆積頂端的最大值換到陣列尾端、堆積縮小一格、把新的頂端往下篩，直到堆積只剩一個元素。

\begin{pseudo}{堆積排序（由小到大）}
\Procedure{Sift-Down}{$A, i, m$} \Comment{堆積是 $A[1..m]$}
    \While{$2i \le m$}
        \State $j \gets 2i$
        \If{$j + 1 \le m$ \AND $A[j + 1] > A[j]$}
            \State $j \gets j + 1$ \Comment{$j$ 是比較大的子節點}
        \EndIf
        \If{$A[i] \ge A[j]$}
            \State \Return
        \EndIf
        \State 交換 $A[i]$ 和 $A[j]$；$i \gets j$
    \EndWhile
\EndProcedure
\Statex
\Procedure{Heapsort}{$A, n$}
    \For{$i \gets \lfloor n/2 \rfloor$ \DownTo $1$} \Comment{第一階段：建堆積}
        \State \Call{Sift-Down}{$A, i, n$}
    \EndFor
    \For{$m \gets n$ \DownTo $2$} \Comment{第二階段}
        \State 交換 $A[1]$ 和 $A[m]$ \Comment{目前的最大值放到 $A[m]$}
        \State \Call{Sift-Down}{$A, 1, m - 1$}
    \EndFor
\EndProcedure
\end{pseudo}

\textbf{正確性}。 \proc{Sift-Down}$(A, i, m)$ 的前提是 $i$ 的兩棵子樹都已經是堆積；它把 $A[i]$ 和較大的子節點交換、往下走，
結束時以 $i$ 為根的子樹成為堆積。 第一階段由下往上做：$\lfloor n/2 \rfloor + 1$ 以後都是葉子，本來就是堆積；處理到 $i$ 時，
它的子節點編號比較大，已經處理過，前提成立，做完 $i = 1$ 整個陣列就是堆積。 第二階段的迴圈不變量：每一輪開始時，
$A[1..m]$ 是最大堆積、裝著最小的 $m$ 個數，$A[m+1..n]$ 是最大的 $n - m$ 個數、已排好。 把 $A[1]$（$A[1..m]$ 的最大值）
換到 $A[m]$，再把 $A[1..m-1]$ 調回堆積（兩棵子樹沒動，仍是堆積），不變量保持；$m = 1$ 時全部排好。

\textbf{時間：第一階段 $O(n)$}。 \proc{Sift-Down} 從高度 $h$ 的節點開始，最多往下走 $h$ 層，每層 $2$ 次比較，成本 $O(h)$。
$n$ 個節點的堆積高度是 $\lfloor \lg n \rfloor$，高度為 $h$ 的節點最多 $\lceil n / 2^{h+1} \rceil$ 個（資結書定理~8.1），所以
\[
\sum_{h=0}^{\lfloor \lg n \rfloor} \left\lceil \frac{n}{2^{h+1}} \right\rceil \cdot O(h)
= O\!\left(n \sum_{h=0}^{\infty} \frac{h}{2^{h}}\right) = O(2n) = O(n),
\]
其中 $\sum_{h \ge 0} h/2^h = 2$（對 $\sum_{h \ge 0} y^h = \frac{1}{1-y}$ 微分再乘 $y$，得 $\sum h y^h = \frac{y}{(1-y)^2}$，代 $y = 1/2$）。
直覺：一半的節點是葉子、根本不用動，越高的節點越少。

\textbf{時間：第二階段 $O(n \log n)$}。 第 $m$ 輪對大小 $m - 1$ 的堆積做 \proc{Sift-Down}，最多 $\lfloor \lg(m - 1) \rfloor$ 層、
每層 $2$ 次比較，所以比較次數不超過
\[
\sum_{m=2}^{n} 2\lfloor \lg(m - 1) \rfloor \le 2\sum_{m=1}^{n-1} \lg m = 2\lg\bigl((n-1)!\bigr) \le 2n \lg n .
\]
合起來 $O(n) + O(n \log n) = O(n \log n)$，而且這是對任何輸入都成立的上界。

\textbf{下界}。 堆積排序是比較式排序，最壞情況至少要 $\lg(n!) = \Theta(n \log n)$ 次比較（資結書第~15.1~節），所以最壞情況是
$\Theta(n \log n)$。 另外，鍵值全部相同時 \proc{Sift-Down} 第一次比較就停，總時間只有 $\Theta(n)$；這說明上面的分析是上界，最壞情況才一定是
$\Theta(n \log n)$。

\textbf{其他性質}。 只用常數個額外變數，是原地排序（額外空間 $O(1)$）；交換會跨越很遠的距離，不穩定。 和快速排序比，堆積排序
保證最壞 $O(n \log n)$；和合併排序比，不需要 $\Theta(n)$ 的額外空間。
\end{pastsoln}
